
\section{Calibrated UAS: Empirical Closed-Form Extension}
\label{sec:calibrated_uas}

The UAS routing law (Theorem~\ref{thm:uas}) is provably stable, but its parameterization --- the energy $\uasenergy{i} = (\Qlen{i}+1)/\srvrate_i$ and prior $\uasprior{i} = \srvrate_i/\totcap$ --- was designed for analytical tractability rather than empirical optimality.
This section introduces Calibrated UAS, a three-parameter generalization that relaxes UAS's structural constraints to improve steady-state performance.

\begin{remark}[Theorem status]
\label{rem:caluas_status}
Calibrated UAS does \textbf{not} carry a formal stability guarantee.
The stability guarantees of Theorem~\ref{thm:uas} apply only to the UAS special case $(\calbeta, \calgamma, \caloffset) = (1, 1, 1)$.
The weighted variational identity (Proposition~\ref{prop:uas_variational}) and Jensen-based arrival bound (Lemma~\ref{lem:uas_arrival}) do not extend to the generalized parameterization.
Establishing stability for general $(\calbeta, \calgamma, \caloffset)$ is an open problem (see Section~\ref{sec:discussion}).
\end{remark}

\begin{definition}[Calibrated UAS]
\label{def:caluas_family}
The Calibrated UAS family is parameterized by $(\calbeta, \calgamma, \caloffset) \in (0,1] \times [0,1] \times [0,1]$:
\begin{equation}
\label{eq:caluas_law}
\rprob{i}^{\caluas}(\Qstate, \srvrate) \propto \srvrate_i^{\calgamma} \exp\!\left(-\temp \frac{\Qlen{i} + \caloffset}{\srvrate_i^{\calbeta}}\right).
\end{equation}
Setting $(\calbeta, \calgamma, \caloffset) = (1, 1, 1)$ recovers UAS~\eqref{eq:uas_law} exactly.
\end{definition}

\subsection{Parameter Semantics}
Each parameter modifies a distinct structural aspect of UAS:
\begin{itemize}
\item $\calbeta < 1$ \emph{tempers} the service-rate normalization: $\srvrate_i^{\calbeta}$ grows more slowly than $\srvrate_i$, reducing the penalty gap between fast and slow servers.
\item $\calgamma < 1$ \emph{softens} the service-rate prior: $\srvrate_i^{\calgamma}$ flattens the routing bias toward fast servers.
\item $\caloffset < 1$ \emph{reduces} the look-ahead offset from $\Qlen{i}+1$ to $\Qlen{i}+\caloffset$, attenuating the anticipatory correction.
\end{itemize}

\subsection{Benchmark Parameters}
The benchmark values used throughout this paper are
\begin{equation}
\label{eq:caluas_params}
(\calbeta, \calgamma, \caloffset) = (0.85,\; 0.5,\; 0.5),
\end{equation}
selected via a grid search over $\calbeta \in \{0.5, 0.7, 0.85, 1.0\}$, $\calgamma \in \{0.25, 0.5, 0.75, 1.0\}$, and $\caloffset \in \{0.25, 0.5, 0.75, 1.0\}$ (64 configurations).
Each configuration was evaluated over 16 independent replications on a validation set separate from the 32-replication anchor benchmark.
The selected parameters minimize mean total queue length on the validation set.
We acknowledge that this grid is coarse and that the parameter surface may not be convex; a more systematic optimization procedure (e.g., Bayesian optimization over a continuous domain) is left for future work.

\subsection{Anchor Benchmark Result}
Calibrated UAS achieves the lowest mean total queue length among all closed-form policies in the anchor benchmark (Section~\ref{sec:exp:policy}):
\begin{center}
\begin{tabular}{@{}lr@{}}
\toprule
\textbf{Policy} & \textbf{Mean $\E[\sum_i\Qlen{i}]$} \\
\midrule
Calibrated UAS & 10.042 \\
JSSQ & 11.016 \\
UAS ($\temp{=}10$) & 11.495 \\
\bottomrule
\end{tabular}
\end{center}



